\section{System 2: The Geometric Impedance}
\label{sec:system2_geometric_impedance}

\CatchFileBetweenTags{\AlphaInvCAPVal}{constants.tex}{AlphaInvCAPVal}
\CatchFileBetweenTags{\AlphaInvMAPVal}{constants.tex}{AlphaInvMAPVal}
\CatchFileBetweenTags{\AlphaInvPROVal}{constants.tex}{AlphaInvPROVal}
\CatchFileBetweenTags{\AlphaInvGOVVal}{constants.tex}{AlphaInvGOVVal}
\CatchFileBetweenTags{\AlphaInvTOLVal}{constants.tex}{AlphaInvTOLVal}
\CatchFileBetweenTags{\AlphaInvMARVal}{constants.tex}{AlphaInvMARVal}
\CatchFileBetweenTags{\AlphaInvVal}{constants.tex}{AlphaInvVal}

Having isolated the unique substrate, we must test whether it can sustain localized persistent structures. A topological defect concentrates information flux against the homogeneous lattice; without an entropic cost accounting, the substrate would dissolve any localized distinction.  We
therefore quantify the minimum geometric action required to maintain one unit of topological charge ($Q_{top}=\chi/2=1$) through a single fundamental cycle ($\tau=1$).

Because the boundary $\chi=2$ enforces a strict binary partition of information flux, charge is a topological invariant.  The \textbf{geometric impedance} is the resulting dimensionless entropic action per unit charge:
\begin{equation}
    Z_{geo} = \frac{S_{\Phi,\text{defect}}}{Q_{top}}
\end{equation}

At quiescent equilibrium the six pillars operate on orthogonal domains, so $Z_{geo}$ is the exact additive sum of their individual geometric costs (\cref{sec:entropic_balance_sheet}).

% ============================================================
% CAP must satisfy all of the following:
%
% 1. BLINDING TEST. Every step must be one we would accept without
%    knowing alpha^-1 ~ 137. Any step persuasive only because the
%    total comes out right is numerology and not acceptable.
% 2. DEPENDENCY DIRECTION. Premises may only be System 0 requirements
%    (Closure, Finiteness, Unitarity, Causality, Homogeneity) and
%    System 1 invariants (Delta=43, c=1, Euclidean substrate metric,
%    periodic causal cycle). Nothing from System 2 onward and no
%    companion-paper dynamics (Wilson loops, gravity strain,
%    Zitterbewegung) may be premises; they may appear only as labeled
%    forward-consistency remarks.
% 3. FORWARD CONSISTENCY. The result may not contradict later
%    structure (Wilson perimeter law, capacity-strain gravity, mass
%    as closed chiral circulation), but consistency is a CHECK, never
%    a premise.
% 4. PI MUST BE CONSTRUCTED. Pi may enter only as the measured
%    invariant of an explicitly constructed geometric object. The
%    microscopic lattice has no continuous rotational symmetry; the
%    isotropic round ball is the macroscopic idealization forced by
%    the Homogeneity requirement (macroscopic observables must be
%    invariant under lattice orientation). The step where isotropy
%    replaces the discrete lattice ball with the round ball must be
%    explicit, named, and scoped to macroscopic entropic cost.
% 5. EVERY FACTOR INTERROGATED. The text must answer, in-line:
%    why the factor is exactly 1; why diameter Delta and not radius
%    Delta (why not 2*pi*Delta); why the perimeter and not the area
%    (why not ~pi*Delta^2/4) or the path length (why not Delta); why
%    the sign is structural (+) under the Entropic Balance Sheet.
% 6. OPERATIONAL UNITS. The result must be a dimensionless count of
%    elementary link actions per fundamental cycle, consistent with
%    S_Phi as a transition count.
% 7. ONE-SENTENCE TEST. The geometric picture must be stateable in
%    one sentence: a persistent defect is a closed circulation whose
%    parts must remain in causal contact within one cycle; its
%    boundary must be re-traced every cycle.
% ============================================================
\subsection{The Resonant Circumference (Capacity)}
\label{subsec:alpha_cap}
Capacity measures the entropic footprint required for a topological defect to maintain structural coherence across its fundamental cycle. The dominant contribution to the vacuum impedance ($\approx 99.9\%$) originates from this fundamental geometry (this corresponds in the standard effective field theory dual to the Wilson Loop). Without this cost, the state would propagate openly at the bandwidth limit ($c=1$) and delocalize completely.

% Establish the metric
At Quiescent Equilibrium, the substrate's footprint is evaluated on its native positive-definite (Euclidean) metric required by Unitarity (\Cref{subsec:positive_definite}), with causal propagation restricted to a $(1+1)D$ plane $(x,\tau)$.

% Derive the maximal reach from periodicity
The fundamental cycle is periodic: the defect's configuration must recur identically after $\Delta$ updates or it does not persist. For any closed trajectory, reaching a distance $d$ from the anchor costs at least $d$ steps outbound and $d$ steps to return, so $2d \le \Delta$: the maximal one-way reach is $d_{\max} = \Delta/2$, saturated by the out-and-back path.

% Construct the invariant envelope
% Bridge discrete tocontinuous
By Homogeneity, the substrate admits no preferred direction. The invariant envelope (the union of all valid closed $\Delta$-cycles anchored at the node) is therefore an isotropic disk of radius $\Delta/2$ in the embedding Euclidean metric, making its diameter ($\Delta/2 + \Delta/2 = \Delta$) exactly match the temporal period. While the lattice provides a discrete address space, the geometric measure of this isotropic envelope on the embedding space is continuous. Any discrete anisotropy cancels in the Homogeneity-enforced average over all orientations, leaving the continuous Euclidean measure.

% Identify the boundary is the cost
% Invoke isotropy to pi
Because the defect is defined strictly by its topological boundary (\Cref{sec:topological_boundary}), and its interior being indistinguishable from the vacuum, the entropic cost is precisely the measure of this continuous perimeter. Given the established isotropic limit, this measure is exactly the circumference of a Euclidean circle, yielding the strict geometric ratio $\frac{C}{\Delta} = \pi$.

\begin{equation}
    Z_{CAP} = C = \pi\Delta \approx \num{\AlphaInvCAPVal}
\end{equation}


\subsection{Topological Boundary (Map)}
\label{subsec:alpha_map}
The Map is the boundary constraint preventing internal state data from dissolving into the background substrate. Without this topological cost, the substrate could not distinguish the system's inside from its outside, causing information to leak into the unquantized background.

As established in \cref{sec:topological_boundary} by the combinatorial Gauss-Bonnet theorem (where total discrete curvature equates to $2\pi\chi$), any closed, regular surface must satisfy a topological boundary condition. To prevent boundaries with holes ($g \ge 1$), where non-contractible loops destroy the strict binary partition of `Inside' versus `Outside' and permit unquantized information flux, the boundary must be simply-connected. The unique simply-connected, closed 2D surface is the sphere ($S^2$), which is characterized by the unique maximum Euler characteristic:
\begin{equation}
    Z_{MAP} = +\chi = \AlphaInvMAPVal
\end{equation}
Because $\chi$, the Euler characteristic, is a pure topological invariant, it is an exact, indivisible integer. It is structurally forbidden from scaling, running, or acquiring continuous corrections.


\subsection{Alignment Efficiency (Protocol)}
\label{subsec:alpha_pro}
The Protocol is the alignment cost mediating phase information across the orthogonal spacetime (Sector A) and internal symmetry (Sector B) channels. Without it, these channels decouple, fragmenting the vacuum into isolated subspaces.

The manifold resolution $R_M = D\Delta$ (\Cref{subsec:manifold_resolution}) is the total available coordination capacity, the maximum distinct spacetime addresses per fundamental cycle. A localized defect consumes exactly $\sigma = 5$ degrees of freedom to define its geometric embedding, leaving residual capacity:

\begin{equation}
    C_{res} = R_M - \sigma = (4 \cdot 43) - 5 = 167
\end{equation}

Entropic impedance is the inverse of flow capacity. The greater the residual capacity, the lower the resistance to phase transport, and because efficient coordination reduces net dissipation, this contribution is an organizational gain, strictly negative per the Entropic Balance Sheet (\Cref{sec:entropic_balance_sheet}):
\begin{equation}
    Z_{PRO} = -\frac{1}{C_{res}} = -\frac{1}{167} \approx \num{\AlphaInvPROVal}
\end{equation}
The linear form is mandated by statelessness: the Protocol coordinates instantaneously, with no recursive memory that would permit higher-order couplings.


\subsection{Stabilizing Potential (Governor)}
\label{subsec:alpha_gov}
The Governor is the stabilizing constraint preventing runaway entropic divergence across the temporal cycle. Without it, local phase excitations undergo unbounded feedback, overflowing node capacity.
 
To satisfy the Homogeneity constraint, this boundary constraint cannot be anchored to any preferred temporal coordinate; it must be distributed maximally uniformly across the fundamental causal loop of length $\Delta = 43$. 

Because the loop is closed and periodic, it possesses no endpoint boundaries (no ``fencepost'' corrections). To prevent non-linear self-interactions which would violate the vacuum's minimal interaction rule (the Toll constraint), the geometric coupling must be strictly first-order (linear). The exact stabilizing potential density per step is therefore the linear ratio of the boundary ($\chi = 2$) to the loop length ($\Delta = 43$). As an organizational gain that prevents entropic divergence, its contribution is negative:
\begin{equation}
    Z_{GOV} = -\frac{\chi}{\Delta} = -\frac{2}{43} \approx \num{\AlphaInvGOVVal}
\end{equation}


\subsection{Entropic Transition (Toll)}
\label{subsec:alpha_tol}
The Toll is the entropic cost of executing a state transition; without it, causal history collapses. Geometrically, it is the probability of selecting a valid transition from the unconstrained phase space:

\begin{equation}
    Z_{TOL} = \frac{\Omega_{valid}}{\Omega_{total}}.
\end{equation}

The numerator, the minimal state transition, is defined by three geometric filters required to prevent non-unitary mixing, unquantized flux, or coordinate self-aliasing:
\begin{enumerate}
    \item \textbf{State Identity (1):} The interaction acts as the identity on internal channel configuration.
    \item \textbf{Topological Flux ($\chi=2$):} The boundary permits exactly two independent flux orientations.
    \item \textbf{Coordinate Non-Degeneracy ($R_M - \sigma$):} The transition target cannot alias with the interaction's own geometric embedding, leaving the orthogonal residual $(R_M - \sigma)$ valid coordinates.
\end{enumerate}

These three filters are independently necessary and jointly sufficient; no fourth filter is available because the $E_8\to D_4\oplus D_4$ projection exhausts the node's orthogonal geometric generators. Because identity (Sector B), topological flux (Sector A), and spatial addresses ($\chi$) govern strictly orthogonal configuration spaces, their independent valid state counts combine multiplicatively. The minimal valid phase space is therefore exactly:

\begin{equation}
    \Omega_{valid} = 1 \cdot \chi \cdot (R_M - \sigma).
\end{equation}

The denominator counts independent dimensions of the unconstrained configuration manifold. On the substrate, a non-trivial state transition requires an interaction; a 2-point operation is mere propagation (identity). By Schur's lemma, invariant couplings of distinct channel types are strictly trilinear; the minimal interaction vertex is the Spin(8) triality singlet $\mathbf{8}_v \otimes \mathbf{8}_s \otimes \mathbf{8}_c \to \mathbf{1}$. Because triality provides exactly three independent 8-dimensional representations, a fourth leg must duplicate a channel type, producing Causal Aliasing. Aliasing collapses distinct states into a single temporal slot, rendering the causal history non-injective and violating Unitarity. Higher-point transitions are composite sequences of 3-point vertices.

With no external observer to assign internal channel roles prior to interaction, an unconstrained 3-point vertex cannot pre-distinguish between Sector A and Sector B degrees of freedom. Rather than being restricted to the single-sector internal bit-depth $\nu=16$, each leg independently samples the total topological node capacity $N=2\nu=32$ across both orthogonal sectors.

The unconstrained phase space of admissible state assignments for a 3-leg vertex is therefore $N^3$. Combining this with the local interaction degrees of freedom ($\sigma=5$) and the available manifold addresses ($R_M=172$), the total unconstrained phase space is the Cartesian product:

\begin{equation}
    \Omega_{total} = N^3 \cdot \sigma \cdot R_M.
\end{equation}

The Toll is therefore the ratio of minimal valid to maximal unconstrained phase space:
\begin{equation}
\begin{split}
    Z_{TOL} &= \frac{\Omega_{valid}}{\Omega_{total}} 
    = \frac{\chi(R_M - \sigma)}{N^3 \sigma R_M} \\
    &\approx \num{\AlphaInvTOLVal}
\end{split}
\end{equation}


\subsection{Resolution Floor (Margin)}
\label{subsec:alpha_mar}
At the substrate level, the Margin manifests as an resolution threshold: the minimum signal amplitude required to be resolved above the geometric noise floor. Below this floor, structured information dissolves into ambient substrate fluctuations.

The minimum resolvable signal is exactly 1 quantum of information diluted across the \emph{minimal} number of distinct configurations ($V_{config}$) required to specify a localized defect. Because the $E_8 \to D_4 \oplus D_4$ projection partitions the lattice into strictly orthogonal sectors, independent specification domains combine multiplicatively (Cartesian product).

The minimal configuration volume requires three domains to exhaust the node's orthogonal geometric generators and prevent underspecification:
\begin{enumerate}
    \item \textbf{The State Vector ($L_{embed} = 31$):} The total description length of a dynamic defect, structurally derived in \cref{subsec:lembed}.
    \item \textbf{The Local Interaction Volume ($V_{local} = \sigma + 1 = 6$):} A localized lattice interaction is a connected subgraph spanning $\sigma = 5$ independent directions from a central node. To strictly localize the interaction to a single neighborhood, the maximum distance from the center to any node must be 1 (graph radius $r=1$). The Star Graph ($S_\sigma$) is the strictly unique connected tree topology of $\sigma$ edges that possesses a radius of 1; any other connected tree has radius $r \ge 2$, delocalizing the interaction across multiple neighborhoods and violating the Orthogonality Condition. The minimal local volume is therefore $\sigma+1=6$.
    \item \textbf{The Causal Area} ($A_{causal} = \Delta^2$): As established in \cref{sec:fundamental_resonance_filter1}, fundamental causality restricts information propagation on the discrete substrate to a local $(1+1)$D plane. Operating at the maximum speed limit $c \equiv 1$, a fundamental cycle of temporal duration $\Delta$ sweeps a 1D spatial horizon of length $\Delta$. The resulting spacetime resolution count is strictly $\Delta \cdot \Delta = \Delta^2$.
\end{enumerate}

Because state vector, local topology, and causal horizon constitute orthogonal configuration spaces, their independent specification domains combine via Cartesian product, yielding a multiplicative volume:
\begin{equation}
    V_{config} = L_{embed} \cdot (\sigma + 1) \cdot \Delta^2
\end{equation}

The Margin impedance is the reciprocal: the probability that a single quantum of structured information remains distinguishable from the geometric noise floor after being diluted across all independent specification domains:
\begin{equation}
\begin{split}
    Z_{MAR} &= \frac{1}{V_{config}} = \frac{1}{L_{embed} \cdot (\sigma + 1) \cdot \Delta^2} \\
    &=\frac{1}{31 \cdot (5+1) \cdot 43^2} \approx \num{\AlphaInvMARVal}
\end{split}
\end{equation}


\subsection{The Additive Baseline and the Finite Sum}
\label{sec:additive_baseline}

% JOB 1+2: Factorization and its justification
The quiescent-equilibrium state sum factorizes exactly over the six pillars:
\begin{equation}
\mathcal{Z}_{\text{total}} = \prod_{i=1}^6 \mathcal{Z}_i,
\end{equation}
because, as established in the Entropic Balance Sheet (\Cref{sec:entropic_balance_sheet}), the six pillars operate on orthogonal, independent domains. At strict zero flux, no propagating excitations exist to mediate interactions between these domains; mutual information vanishes identically, forbidding cross-terms.

% JOB 3+4: Unified currency and additive form
As established in the Entropic Balance Sheet, the entropic action is the negative logarithm of the factorized state sum. Each pillar therefore contributes an exact exponential weight $\mathcal{Z}_i = e^{-Z_i}$, where $Z_i$ is the dimensionless geometric action exacted by that pillar. The geometric impedance is the exact additive entropy:
\begin{equation}
    Z_{\text{geo}} = -\ln \mathcal{Z}_{\text{total}} = \sum_{i=1}^6 Z_i.
\end{equation}

% JOB 5: Heterogeneity defense
Despite differing geometric origins (boundary measure, topological invariant, or phase-space ratio), each term is a dimensionless ratio relative to the fundamental lattice unit cell, rendering all terms commensurable in the dimensionless geometric action $Z_{\text{geo}}$.

% JOB 6+7+8: Completeness, signs, and transition
Two constraints rigidly bound this finite sum, ensuring it forms a complete and minimal basis:
\begin{enumerate}
    \item \textbf{Individual Forcing and Completeness:} The functional form of each term is strictly forced by the geometric and information-theoretic requirements of its component. These six components exhaust the architectural degrees of freedom required for persistence. The structural invariants $\mathbb{S} = \{\Delta=43, \nu=16, \sigma=5, \chi=2\}$ exhaust the independent invariants of the projection. \textit{(Note that $\nu=16$ enters only through the total node capacity $N=2\nu=32$; no term involves $\nu$ independently because the chiral bit-depth is fully allocated to the channel structure).}
    
    \item \textbf{Sign Classification:} The signs are fixed prior to evaluation by the Entropic Balance Sheet (\Cref{sec:entropic_balance_sheet}), confirmed by the failure mode analysis. Terms representing irreducible structural overhead (Capacity, Map, Toll, Margin) are strictly positive. Terms representing organizational gains (Protocol, Governor) act analogously to negative feedback in control theory, reducing net dissipation, and are therefore strictly negative.
\end{enumerate}

The geometric impedance is therefore exactly evaluated by expanding this sum into the individual geometric costs derived in \crefrange{subsec:alpha_cap}{subsec:alpha_mar}:

\begin{equation}\label{eq:alpha_inverse}
\begin{split}
Z_{\text{geo}} = \underbrace{\pi\Delta}_{\text{CAP}}
+ \,\underbrace{\chi}_{\text{MAP}}
- \,\underbrace{\frac{1}{R_M - \sigma}}_{\text{PRO}}
- \,\underbrace{\frac{\chi}{\Delta}}_{\text{GOV}} \\
+ \,\underbrace{\frac{1 \cdot \chi \cdot (R_M - \sigma)}{N^3 \cdot \sigma \cdot R_M}}_{\text{TOL}}
+ \,\underbrace{\frac{1}{L_{\text{embed}} \cdot (\sigma + 1) \cdot \Delta^2}}_{\text{MAR}}
\end{split}
\end{equation}


\subsection{Numerical Evaluation}
\label{subsec:alpha_numerical_validation}
Summing the geometric components:
\begin{equation}
    Z_{\text{geo}} = \num{\AlphaInvVal}\dots
\end{equation}

Every term is an exact rational combination of the structural invariants except $\pi$, which enters as the universal isotropic ratio$C/\Delta$ (Section~\ref{subsec:alpha_cap}). The derivation contains zero continuous empirical parameters; precision is therefore limited only by the evaluation of $\pi$.